\documentclass[aspectratio=169,11pt]{beamer}

% =========================================================
% TEMPLATE MAESTRO
% =========================================================
\input{~/Dropbox/Latex/TemplateSlides/template-slide.tex}

% =========================================================
% INFORMACIÓN INSTITUCIONAL
% =========================================================

\renewcommand{\CourseName}{Diseño Web 1}
\renewcommand{\CourseCode}{ADSO-CTMA}
\renewcommand{\Faculty}{Análisis y Desarrollo de Software}
\renewcommand{\Institution}{Servicio Nacional de Aprendizaje SENA}
\renewcommand{\Instructor}{Deimer Miranda Montoya, MSc.(c).}
\renewcommand{\Email}{damiranda@sena.edu.co}
\renewcommand{\InstitutionLogo}{Template/LogoSena.png}

\begin{document}

    \TitleSlide
    \NormasSlide

    \SectionSlide
    {Diseño web responsivo}
    {Interfaces que se adaptan al espacio disponible}

    % =========================================================
    \section{Punto de partida}
    % =========================================================

    \begin{frame}{¿Qué aprenderemos?}
        Al finalizar la clase podrá:

        \begin{itemize}
            \item Explicar por qué una página debe adaptarse a distintos anchos.
            \item Configurar correctamente el \textit{viewport}.
            \item Construir distribuciones fluidas.
            \item Adaptar imágenes y contenido multimedia.
            \item Crear media queries a partir del contenido.
            \item Aplicar un enfoque \textit{mobile first}.
            \item Combinar diseño responsivo con Flexbox y Grid.
            \item Verificar el resultado mediante las herramientas del navegador.
        \end{itemize}
    \end{frame}

    \begin{frame}{Una página no tiene un único tamaño}
        El mismo sitio puede visitarse desde:

        \begin{itemize}
            \item un teléfono sostenido verticalmente;
            \item una tableta;
            \item un computador portátil;
            \item una pantalla de escritorio amplia;
            \item una ventana que ocupa solo una parte del monitor.
        \end{itemize}

        \Idea{
            No diseñamos para una lista cerrada de dispositivos.
            Diseñamos para el espacio disponible.
        }
    \end{frame}

    \begin{frame}{El problema de un diseño rígido}
        \begin{center}
            \begin{tikzpicture}[
                pantalla/.style={draw=BrandBlue, very thick, rounded corners},
                bloque/.style={draw=PracticeTitle, fill=PracticeBody, thick,
                    rounded corners=2pt}
                ]
                \draw[pantalla] (0,0) rectangle (6.7,3.6);
                \node[font=\bfseries] at (3.35,3.95) {Pantalla amplia};
                \draw[bloque] (0.35,2.65) rectangle (6.35,3.25);
                \draw[bloque] (0.35,0.45) rectangle (2.1,2.35);
                \draw[bloque] (2.35,0.45) rectangle (6.35,2.35);

                \draw[pantalla] (8.5,0) rectangle (10.7,3.6);
                \node[font=\bfseries] at (9.6,3.95) {Pantalla estrecha};
                \draw[bloque] (8.8,2.65) rectangle (14.8,3.25);
                \draw[bloque] (8.8,0.45) rectangle (10.55,2.35);
                \draw[bloque] (10.8,0.45) rectangle (14.8,2.35);
                \draw[->, very thick, WarnTitle] (10.9,1.3) -- (12.1,1.3)
                node[right] {desbordamiento};
            \end{tikzpicture}
        \end{center}

        \Warning{
            Si todos los tamaños son fijos, el contenido puede salir de la pantalla
            y obligar al usuario a desplazarse horizontalmente.
        }
    \end{frame}

    \begin{frame}{¿Qué significa responsivo?}
        \DefBlock[Diseño web responsivo]{
            Es una estrategia de diseño y desarrollo que permite que una interfaz
            reorganice, redimensione o simplifique su contenido según el espacio disponible.
        }

        No consiste únicamente en ``hacer todo más pequeño''.

        \begin{itemize}
            \item El contenido conserva su legibilidad.
            \item La estructura puede cambiar.
            \item Los controles siguen siendo utilizables.
            \item La información importante mantiene su prioridad.
        \end{itemize}
    \end{frame}

    \begin{frame}{Tres bases del diseño responsivo}
        \begin{center}
            \begin{tikzpicture}[
                base/.style={draw=BrandBlue, thick, rounded corners,
                    minimum width=3.2cm, minimum height=1.2cm, align=center}
                ]
                \node[base, fill=IdeaBody] at (1.8,2) {Distribución\\flexible};
                \node[base, fill=PracticeBody] at (6,2) {Imágenes y medios\\adaptables};
                \node[base, fill=ResultBody] at (10.2,2) {Reglas según\\el espacio};
                \draw[->, >=Latex, thick] (3.45,2) -- (4.35,2);
                \draw[->, >=Latex, thick] (7.65,2) -- (8.55,2);
            \end{tikzpicture}
        \end{center}

        \Result{
            Las medidas relativas aportan fluidez; las media queries permiten
            cambiar la composición cuando esa fluidez deja de ser suficiente.
        }
    \end{frame}

    % =========================================================
    \section{Proyecto de prueba}
    % =========================================================

    \SectionSlide
    {Proyecto de prueba}
    {Una reserva natural que crecerá paso a paso}

    \begin{frame}[fragile]{Estructura inicial}
        Cree la siguiente estructura:

        \begin{bashcodedark}
            reserva-andina/
            ├── index.html
            ├── css/
            │   └── styles.css
            └── img/
                ├── portada.jpg
                ├── sendero.jpg
                └── aves.jpg
        \end{bashcodedark}

        \Practice{
            Utilice imágenes propias o libres. Los nombres pueden cambiar,
            siempre que las rutas del HTML también se actualicen.
        }
    \end{frame}

    \begin{frame}[fragile]{HTML básico: encabezado}
        \begin{htmlcodedark}
        <!DOCTYPE html>
        <html lang="es">
            <head>
                <meta charset="UTF-8">
                <meta name="viewport"
                content="width=device-width, initial-scale=1.0">
                <title>Reserva Andina</title>
                <link rel="stylesheet" href="css/styles.css">
            </head>
            <body>
                <header class="encabezado">
                    <h1>Reserva Andina</h1>
                    <p>Naturaleza, senderos y conservación</p>
                </header>
        \end{htmlcodedark}
    \end{frame}

    \begin{frame}[fragile]{HTML básico: navegación y portada}
        Continúe dentro de \texttt{body}:

        \begin{htmlcodedark}
            <nav class="navegacion">
                <a href="#inicio">Inicio</a>
                <a href="#rutas">Rutas</a>
                <a href="#visita">Visita</a>
            </nav>

            <main id="inicio">
                <section class="portada">
                    <img src="img/portada.jpg"
                        alt="Montañas de la reserva">
                    <div>
                        <h2>Explore el bosque altoandino</h2>
                        <p>Planee una visita responsable.</p>
                    </div>
                </section>
        \end{htmlcodedark}
    \end{frame}

    \begin{frame}[fragile]{HTML básico: tarjetas}
        \begin{htmlcodedark}
            <section id="rutas" class="rutas">
                <article class="tarjeta">
                    <img src="img/sendero.jpg" alt="Sendero del agua">
                    <h3>Sendero del agua</h3>
                    <p>Recorrido de dificultad baja.</p>
                </article>

                <article class="tarjeta">
                    <img src="img/aves.jpg" alt="Ave sobre una rama">
                    <h3>Ruta de las aves</h3>
                    <p>Observación guiada al amanecer.</p>
                </article>
            </section>
        </main>
        </body>
    </html>
        \end{htmlcodedark}
    \end{frame}

    \begin{frame}[fragile]{CSS inicial}
        \begin{csscodedark}
            * {
                box-sizing: border-box;
            }

            body {
                margin: 0;
                font-family: Arial, sans-serif;
                color: #18332b;
                background-color: #f4f7f5;
            }

            img {
                display: block;
            }
        \end{csscodedark}

        Primero construiremos una versión funcional para espacios pequeños.
    \end{frame}

    % =========================================================
    \section{Viewport}
    % =========================================================

    \SectionSlide
    {Viewport}
    {El área visible de la página}

    \begin{frame}{La pantalla y el viewport no son lo mismo}
        \begin{center}
            \begin{tikzpicture}[
                marco/.style={draw=BrandBlue, very thick, rounded corners=8pt},
                pagina/.style={draw=PracticeTitle, thick, fill=PracticeBody}
                ]
                \draw[marco] (0,0) rectangle (4.2,5.4);
                \draw[pagina] (0.35,0.55) rectangle (3.85,4.85);
                \node at (2.1,4.45) {contenido visible};
                \draw[<->, >=Latex, thick] (0.35,0.25) -- (3.85,0.25)
               node[midway, below=5mm] {ancho del viewport};

                \node[align=left, text width=6.1cm] at (8.2,3.2) {
                    El \textit{viewport} es la región del navegador donde
                    se representa la página.\\[0.4cm]
                    CSS toma decisiones a partir de ese espacio disponible,
                    no del tamaño físico total del dispositivo.
                };
            \end{tikzpicture}
        \end{center}
    \end{frame}

    \begin{frame}[fragile]{Configurar el viewport}
        Esta etiqueta debe estar dentro de \texttt{head}:

        \begin{htmlcodedark}
            <meta name="viewport"
            content="width=device-width, initial-scale=1.0">
        \end{htmlcodedark}

        \begin{description}
            \item[\texttt{width=device-width}:]
            relaciona el ancho del documento con el viewport del dispositivo.

            \item[\texttt{initial-scale=1.0}:]
            establece la escala inicial normal.
        \end{description}

        \Warning{
            Sin esta configuración, un teléfono puede simular un viewport amplio
            y reducir visualmente toda la página.
        }
    \end{frame}

    \begin{frame}{No bloquee el zoom}
        Evite configuraciones como:

        \begin{center}
            \texttt{user-scalable=no}
        \end{center}

        También conviene evitar límites innecesarios sobre la escala máxima.

        \Idea{
            Algunas personas necesitan ampliar la interfaz para leerla.
            Impedir el zoom crea una barrera de accesibilidad.
        }
    \end{frame}

    % =========================================================
    \section{Distribuciones fluidas}
    % =========================================================

    \SectionSlide
    {Distribuciones fluidas}
    {Tamaños que aprovechan el espacio sin desbordarse}

    \begin{frame}{Tamaño fijo frente a tamaño flexible}
        \begin{center}
            \small
            \begin{tabular}{p{4.3cm}p{5.3cm}}
                \toprule
                \textbf{Tamaño fijo} & \textbf{Tamaño flexible} \\
                \midrule
                \texttt{width: 1000px} & \texttt{width: 90\%} \\
                Conserva el mismo ancho & Depende del contenedor \\
                Puede desbordarse & Puede crecer o reducirse \\
                Útil en casos concretos & Útil para la estructura general \\
                \bottomrule
            \end{tabular}
        \end{center}

        \Compare{
            Las medidas fijas no son incorrectas. El problema aparece cuando
            se aplican sin considerar el espacio realmente disponible.
        }
    \end{frame}

    \begin{frame}[fragile]{Un contenedor flexible con límite}
        \begin{csscodedark}
            main {
                width: min(92%, 1100px);
                margin-inline: auto;
            }
        \end{csscodedark}

        Esta declaración combina dos necesidades:

        \begin{itemize}
            \item en pantallas estrechas utiliza aproximadamente el 92\%;
            \item en pantallas amplias no supera los 1100px;
            \item los márgenes automáticos centran el contenido.
        \end{itemize}

        \Result{
            El contenido puede adaptarse sin formar líneas de texto excesivamente largas.
        }
    \end{frame}

    \begin{frame}[fragile]{Alternativa con width y max-width}
        La misma idea puede expresarse de forma más tradicional:

        \begin{csscodedark}
            main {
                width: 92%;
                max-width: 1100px;
                margin-left: auto;
                margin-right: auto;
            }
        \end{csscodedark}

        \Idea{
            \texttt{max-width} establece un límite, pero permite que el elemento
            sea más estrecho cuando el contenedor no ofrece suficiente espacio.
        }
    \end{frame}

    \begin{frame}{Unidades que conviene reconocer}
        \begin{center}
            \small
            \begin{tabular}{p{2.2cm}p{3.8cm}p{7cm}}
                \toprule
                \textbf{Unidad} & \textbf{Referencia} & \textbf{Uso habitual} \\
                \midrule
                \texttt{px} & unidad CSS & bordes y medidas controladas \\
                \texttt{\%} & contenedor & anchos fluidos \\
                \texttt{rem} & tamaño raíz & texto y espacios consistentes \\
                \texttt{vw} & ancho del viewport & tamaños ligados a la ventana \\
                \texttt{vh} & alto del viewport & secciones relacionadas con la altura \\
                \texttt{fr} & espacio de Grid & distribución de columnas y filas \\
                \bottomrule
            \end{tabular}
        \end{center}
    \end{frame}

    \begin{frame}[fragile]{Espaciado relativo}
        \begin{csscodedark}
            .encabezado,
            .navegacion,
            .portada,
            .rutas {
                padding: 1rem;
            }

            .rutas {
                display: grid;
                gap: 1rem;
            }
        \end{csscodedark}

        \DefBlock[La unidad rem]{
            \texttt{1rem} corresponde al tamaño de fuente del elemento raíz.
            Normalmente equivale a 16px si el usuario no ha cambiado la configuración.
        }
    \end{frame}

    \begin{frame}{Cuidado con 100vw}
        \texttt{width: 100vw} representa todo el ancho del viewport.

        En algunos navegadores puede incluir el espacio de la barra de desplazamiento,
        lo que genera un pequeño desbordamiento horizontal.

        \Warning{
            Para un bloque normal que debe ocupar el ancho de su contenedor,
            generalmente es más seguro utilizar \texttt{width: 100\%}.
        }
    \end{frame}

    % =========================================================
    \section{Imágenes adaptables}
    % =========================================================

    \SectionSlide
    {Imágenes adaptables}
    {Evitar desbordamientos y recortes accidentales}

    \begin{frame}[fragile]{La regla base para las imágenes}
        \begin{csscodedark}
            img {
                display: block;
                max-width: 100%;
                height: auto;
            }
        \end{csscodedark}

        \begin{description}
            \item[\texttt{max-width: 100\%}:]
            impide que la imagen supere el ancho de su contenedor.

            \item[\texttt{height: auto}:]
            conserva la proporción original.

            \item[\texttt{display: block}:]
            elimina el espacio inferior propio de las imágenes en línea.
        \end{description}
    \end{frame}

    \begin{frame}{Conservar la proporción evita deformaciones}
        \begin{center}
            \begin{tikzpicture}[
                foto/.style={draw=BrandBlue, thick, fill=IdeaBody},
                flecha/.style={->, >=Latex, thick}
                ]
                \draw[foto] (0,0.6) rectangle (4,3.1);
                \node at (2,1.85) {imagen original};
                \draw[flecha] (4.4,1.85) -- (5.5,1.85);
                \draw[foto] (5.9,1.05) rectangle (9.1,3.05);
                \node at (7.5,2.05) {proporción};
                \draw[flecha, WarnTitle] (9.5,1.85) -- (10.6,1.85);
                \draw[draw=WarnTitle, thick, fill=WarnBody]
                (11,0.3) rectangle (12.8,3.4);
                \node[rotate=90] at (11.9,1.85) {deformada};
            \end{tikzpicture}
        \end{center}

        \Idea{
            Si se fija el ancho, permita que el navegador calcule el alto;
            si se fija el alto, controle el recorte conscientemente.
        }
    \end{frame}

    \begin{frame}[fragile]{Imágenes uniformes en tarjetas}
        Cuando todas las tarjetas deben conservar la misma altura visual:

        \begin{csscodedark}
            .tarjeta img {
                width: 100%;
                height: 220px;
                object-fit: cover;
            }
        \end{csscodedark}

        \DefBlock[object-fit: cover]{
            La imagen cubre el área asignada sin deformarse. Si las proporciones
            no coinciden, una parte de la imagen puede recortarse.
        }

        Use un texto alternativo que describa el contenido importante de la imagen.
    \end{frame}

    \begin{frame}{cover frente a contain}
        \begin{center}
\begin{tikzpicture}[
    marco/.style={
        draw=BrandBlue,
        very thick,
        fill=gray!5
    },
    imagen/.style={
        draw=PracticeTitle,
        thick,
        fill=PracticeBody
    }
    ]

    % =====================================================
    % OBJECT-FIT: COVER
    % =====================================================

    \node[font=\bfseries] at (2.6,3.7) {cover};

    % Parte visible de la imagen
    \draw[imagen]
    (0,0.7) rectangle (5.2,3.2);

    % Tamaño completo que tendría la imagen
    \draw[
    PracticeTitle,
    thick,
    dashed
    ]
    (-0.65,0.7) rectangle (5.85,3.2);

    % Contenedor
    \draw[marco, fill=none]
    (0,0.7) rectangle (5.2,3.2);

    \node[align=center] at (2.6,1.95)
    {cubre todo el espacio\\y recorta los extremos};

    \node[
    font=\scriptsize,
    text=WarnTitle
    ] at (2.6,0.3)
    {parte de la imagen queda fuera};

   % =====================================================
   % OBJECT-FIT: CONTAIN
   % =====================================================

   \begin{scope}[xshift=0.5cm]

       \node[font=\bfseries] at (8.5,3.7) {contain};

       % Contenedor
       \draw[marco]
       (5.9,0.7) rectangle (11.1,3.2);

       % Imagen completa dentro del contenedor
       \draw[imagen]
       (6.7,0.7) rectangle (10.3,3.2);

       \node[align=center] at (8.5,1.95)
       {muestra toda la imagen\\
           y deja espacio libre};

   \end{scope}

\end{tikzpicture}
        \end{center}
    \end{frame}

    % =========================================================
    \section{Media queries}
    % =========================================================

    \SectionSlide
    {Media queries}
    {Aplicar reglas cuando cambia el espacio disponible}

    \begin{frame}{¿Cuándo necesitamos una media query?}
        Primero permita que el diseño sea flexible.

        Después reduzca o amplíe la ventana y observe:

        \begin{itemize}
            \item ¿el texto deja de ser legible?;
            \item ¿las tarjetas quedan demasiado estrechas?;
            \item ¿la navegación ya no cabe?;
            \item ¿se desperdicia demasiado espacio?;
            \item ¿el orden visual sigue siendo claro?
        \end{itemize}

        \Idea{
            Una media query se justifica cuando el contenido necesita una nueva composición.
        }
    \end{frame}

    \begin{frame}[fragile]{Estructura de una media query}
        \begin{csscodedark}
            @media (min-width: 700px) {
                .rutas {
                    grid-template-columns: repeat(2, 1fr);
                }
            }
        \end{csscodedark}

        \begin{description}
            \item[\texttt{@media}:] inicia la consulta.
            \item[\texttt{min-width: 700px}:] expresa la condición.
            \item[Las llaves externas:] contienen las reglas condicionadas.
            \item[La regla interna:] solo se aplica cuando se cumple la condición.
        \end{description}
    \end{frame}

    \begin{frame}{Leer min-width correctamente}
        \begin{center}
            \begin{tikzpicture}[x=0.012cm]
                \draw[very thick, ->, >=Latex] (320,0) -- (1200,0);
                \foreach \x/\t in {360/360,700/700,1000/1000}{
                    \draw[thick] (\x,-0.18) -- (\x,0.18);
                    \node[below] at (\x,-0.22) {\t px};
                }
                \draw[very thick, PracticeTitle] (700,0.45) -- (1180,0.45);
                \node[above, PracticeTitle] at (940,0.45)
                {se aplican las reglas de \texttt{min-width: 700px}};
            \end{tikzpicture}
        \end{center}

        \Result{
            La condición se cumple desde 700px en adelante, incluyendo exactamente 700px.
        }
    \end{frame}

    \begin{frame}[fragile]{Leer max-width correctamente}
        \begin{csscodedark}
            @media (max-width: 699px) {
                .navegacion {
                    flex-direction: column;
                }
            }
        \end{csscodedark}

        La regla se aplica cuando el viewport mide 699px o menos.

        \Compare{
            \texttt{min-width} construye hacia espacios mayores.
            \texttt{max-width} introduce cambios hacia espacios menores.
        }
    \end{frame}

    \begin{frame}{Los breakpoints no son modelos de teléfono}
        \DefBlock[Breakpoint]{
            Es el punto a partir del cual el contenido requiere un cambio de composición.
        }

        Evite pensar únicamente en:

        \begin{center}
            móvil $\rightarrow$ tableta $\rightarrow$ computador
        \end{center}

        \Idea{
            Pruebe el diseño con diferentes anchos y establezca el breakpoint
            donde el contenido comienza a fallar, no donde empieza un dispositivo famoso.
        }
    \end{frame}

    \begin{frame}[fragile]{Combinar condiciones}
        Puede establecer un intervalo:

        \begin{csscodedark}
            @media (min-width: 700px) and (max-width: 999px) {
                .portada {
                    grid-template-columns: 1fr 1fr;
                }
            }
        \end{csscodedark}

        También existen consultas por orientación:

        \begin{csscodedark}
            @media (orientation: landscape) {
                /* Reglas para una vista más ancha que alta */
            }
        \end{csscodedark}

        Use condiciones adicionales solo cuando resuelvan una necesidad real.
    \end{frame}

    \begin{frame}{El orden de las reglas importa}
        Cuando dos reglas tienen la misma especificidad y ambas se cumplen,
        normalmente gana la que aparece después.

        \begin{center}
            \begin{tikzpicture}[
                regla/.style={draw=BrandBlue, thick, rounded corners,
                    minimum width=3.6cm, minimum height=0.85cm}
                ]
                \node[regla, fill=IdeaBody] (base) at (2.2,2.4) {regla base};
                \node[regla, fill=PracticeBody] (media) at (8.2,2.4) {regla en media query};
                \draw[->, >=Latex, very thick] (4.1,2.4) -- (6.0,2.4)
                node[midway, above] {después};
                \node[align=center] at (4.7,0.9)
                {si la condición se cumple,\\la segunda puede sobrescribir la primera};
            \end{tikzpicture}
        \end{center}
    \end{frame}

    % =========================================================
    \section{Mobile first}
    % =========================================================

    \SectionSlide
    {Mobile first}
    {Comenzar con lo esencial y ampliar progresivamente}

    \begin{frame}{¿Qué significa mobile first?}
        \DefBlock[Enfoque mobile first]{
            Las reglas base resuelven primero los espacios estrechos.
            Después se agregan mejoras para anchos mayores mediante \texttt{min-width}.
        }

        No significa que el escritorio sea menos importante.

        \begin{itemize}
            \item Obliga a priorizar el contenido.
            \item Reduce la cantidad de sobrescrituras.
            \item Produce una progresión CSS fácil de seguir.
        \end{itemize}
    \end{frame}

    \begin{frame}{La progresión mobile first}
        \begin{center}
           \begin{tikzpicture}[
               node distance=0.8cm,
               etapa/.style={
                   draw=BrandBlue,
                   thick,
                   rounded corners,
                   text width=3.2cm,
                   minimum height=1.1cm,
                   inner sep=4pt,
                   align=center,
                   font=\small
               },
               flecha/.style={
                   ->,
                   >=Latex,
                   very thick
               }
               ]

               \node[
               etapa,
               fill=IdeaBody
               ] (m) {
                   Reglas base\\
                   una columna
               };

               \node[
               etapa,
               fill=PracticeBody,
               right=of m
               ] (t) {
                   \texttt{min-width: 700px}\\
                   dos columnas
               };

               \node[
               etapa,
               fill=ResultBody,
               right=of t
               ] (d) {
                   \texttt{min-width: 1000px}\\
                   composición amplia
               };

               \draw[flecha] (m) -- (t);
               \draw[flecha] (t) -- (d);

           \end{tikzpicture}
        \end{center}

        Las reglas se acumulan: la vista amplia conserva la base y añade cambios.
    \end{frame}

    \begin{frame}[fragile]{Paso 1: navegación para poco espacio}
        La regla base organiza los enlaces verticalmente:

        \begin{csscodedark}
            .navegacion {
                display: flex;
                flex-direction: column;
                gap: 0.5rem;
                padding: 1rem;
                background-color: #dcebe3;
            }

            .navegacion a {
                padding: 0.75rem;
                color: #18332b;
                text-decoration: none;
            }
        \end{csscodedark}
    \end{frame}

    \begin{frame}[fragile]{Paso 2: navegación horizontal}
        Cuando existe suficiente ancho:

        \begin{csscodedark}
            @media (min-width: 700px) {
                .navegacion {
                    flex-direction: row;
                    justify-content: center;
                }
            }
        \end{csscodedark}

        \Result{
            No fue necesario repetir colores, espacios ni tipografía.
            La media query cambia únicamente la dirección y la distribución.
        }
    \end{frame}

    \begin{frame}[fragile]{Paso 3: portada en una columna}
        Para el ancho menor:

        \begin{csscodedark}
            .portada {
                display: grid;
                gap: 1rem;
                margin-block: 1rem;
                background-color: white;
            }

            .portada > div {
                padding: 1rem;
            }
        \end{csscodedark}

        La imagen aparece primero y el texto después, siguiendo el orden del HTML.
    \end{frame}

    \begin{frame}[fragile]{Paso 4: portada en dos columnas}
        \begin{csscodedark}
            @media (min-width: 700px) {
                .portada {
                    grid-template-columns: 1.2fr 1fr;
                    align-items: center;
                }
            }
        \end{csscodedark}

        \begin{center}
            \begin{tikzpicture}[
                bloque/.style={draw=PracticeTitle, thick, fill=PracticeBody,
                    minimum height=1.3cm}
                ]
                \node[bloque, minimum width=5cm] at (3.2,1.2) {imagen};
                \node[bloque, minimum width=4cm] at (8,1.2) {texto};
            \end{tikzpicture}
        \end{center}
    \end{frame}

    \begin{frame}[fragile]{Paso 5: tarjetas que cambian de columnas}
        Regla base:

        \begin{csscodedark}
            .rutas {
                display: grid;
                grid-template-columns: 1fr;
                gap: 1rem;
            }
        \end{csscodedark}

        Para espacios mayores:

        \begin{csscodedark}
            @media (min-width: 700px) {
                .rutas {
                    grid-template-columns: repeat(2, 1fr);
                }
            }
        \end{csscodedark}
    \end{frame}

    \begin{frame}{Resultado conceptual}
        \begin{center}
            \begin{tikzpicture}[
                scale=0.8,
                transform shape,
                marco/.style={draw=BrandBlue, very thick, rounded corners=5pt},
                bloque/.style={draw=PracticeTitle, fill=PracticeBody, thick,
                    rounded corners=2pt}
                ]
                \node[font=\bfseries] at (2,5.4) {Vista estrecha};
                \draw[marco] (0,0) rectangle (4,5);
                \draw[bloque] (0.3,4.2) rectangle (3.7,4.7);
                \draw[bloque] (0.3,3.35) rectangle (3.7,3.95);
                \draw[bloque] (0.3,2.05) rectangle (3.7,3.1);
                \draw[bloque] (0.3,1.1) rectangle (3.7,1.8);
                \draw[bloque] (0.3,0.25) rectangle (3.7,0.95);

                \node[font=\bfseries] at (9.4,5.4) {Vista amplia};
                \draw[marco] (5.1,0) rectangle (13.7,5);
                \draw[bloque] (5.4,4.2) rectangle (13.4,4.7);
                \draw[bloque] (5.4,3.5) rectangle (13.4,4);
                \draw[bloque] (5.4,1.9) rectangle (10,3.25);
                \draw[bloque] (10.2,1.9) rectangle (13.4,3.25);
                \draw[bloque] (5.4,0.35) rectangle (9.2,1.6);
                \draw[bloque] (9.6,0.35) rectangle (13.4,1.6);
            \end{tikzpicture}
        \end{center}

        El contenido no desaparece: cambia su organización.
    \end{frame}

    \begin{frame}{Mobile first frente a desktop first}
        \begin{center}
            \small
            \begin{tabular}{p{4.5cm}p{4.5cm}}
                \toprule
                \textbf{Mobile first} & \textbf{Desktop first} \\
                \midrule
                Base para poco espacio & Base para espacio amplio \\
                Crece con \texttt{min-width} & Reduce con \texttt{max-width} \\
                Prioriza lo esencial desde el inicio & Puede partir de una composición rica \\
                Suele requerir menos correcciones & Puede necesitar más sobrescrituras \\
                \bottomrule
            \end{tabular}
        \end{center}

        \Idea{
            Ambos enfoques son válidos. En esta clase utilizamos mobile first
            porque facilita una progresión acumulativa.
        }
    \end{frame}

    % =========================================================
    \section{Patrones responsivos}
    % =========================================================

    \SectionSlide
    {Patrones responsivos}
    {Flexbox y Grid como aliados}

    \begin{frame}[fragile]{Patrón Grid sin breakpoint explícito}
        Algunas cuadrículas pueden adaptarse con una sola regla:

        \begin{csscodedark}
            .rutas {
                display: grid;
                grid-template-columns:
                repeat(auto-fit, minmax(min(100%, 260px), 1fr));
                gap: 1rem;
            }
        \end{csscodedark}

        \begin{itemize}
            \item Cada tarjeta intenta conservar una base de 260px.
            \item Nunca supera el ancho disponible.
            \item Grid crea tantas columnas como puedan caber.
        \end{itemize}
    \end{frame}

    \begin{frame}{¿Por qué aparece min(100\%, 260px)?}
        Si el contenedor mide menos de 260px, un mínimo rígido de 260px
        podría producir desbordamiento.

        \begin{center}
            \Large
            \texttt{min(100\%, 260px)}
        \end{center}

        selecciona el valor más pequeño entre:

        \begin{itemize}
            \item todo el ancho disponible;
            \item la base deseada de 260px.
        \end{itemize}

        \Result{La tarjeta cabe incluso dentro de un contenedor muy estrecho.}
    \end{frame}

    \begin{frame}[fragile]{Patrón Flexbox con ajuste de línea}
        \begin{csscodedark}
            .lista-datos {
                display: flex;
                flex-wrap: wrap;
                gap: 1rem;
            }

            .dato {
                flex: 1 1 220px;
            }
        \end{csscodedark}

        Cada elemento parte de 220px, puede crecer y también cambiar de línea.

        \Compare{
            Flexbox funciona muy bien cuando domina una dirección.
            Grid resulta útil cuando las columnas deben permanecer relacionadas.
        }
    \end{frame}

    \begin{frame}{¿Media query o cuadrícula automática?}
        Utilice una cuadrícula automática cuando:

        \begin{itemize}
            \item solo cambia la cantidad de columnas;
            \item todas las tarjetas siguen el mismo patrón;
            \item no es necesario reorganizar otras partes de la página.
        \end{itemize}

        Utilice una media query cuando:

        \begin{itemize}
            \item cambia la composición general;
            \item la navegación necesita otro comportamiento;
            \item se modifican alineación, orden o proporciones;
            \item varios elementos deben cambiar de manera coordinada.
        \end{itemize}
    \end{frame}

    \begin{frame}[fragile]{Contenido que nunca debe desbordarse}
        Palabras extensas, enlaces y fragmentos técnicos pueden romper la composición:

        \begin{csscodedark}
            .tarjeta,
            .contenido {
                min-width: 0;
            }

            .contenido p,
            .contenido a {
                overflow-wrap: anywhere;
            }
        \end{csscodedark}

        \Warning{
            No oculte el problema indiscriminadamente con \texttt{overflow: hidden}.
            El contenido debe seguir disponible para el usuario.
        }
    \end{frame}

    % =========================================================
    \section{Tipografía y controles}
    % =========================================================

    \SectionSlide
    {Tipografía y controles}
    {Leer y utilizar la interfaz en cualquier tamaño}

    \begin{frame}[fragile]{Tipografía fluida con clamp}
        \begin{csscodedark}
            h1 {
                font-size: clamp(2rem, 5vw, 4rem);
            }
        \end{csscodedark}

        \texttt{clamp()} recibe tres valores:

        \begin{center}
            \texttt{clamp(mínimo, preferido, máximo)}
        \end{center}

        \begin{itemize}
            \item Nunca será menor de 2rem.
            \item Puede crecer con el viewport.
            \item Nunca será mayor de 4rem.
        \end{itemize}
    \end{frame}

    \begin{frame}{No reduzca demasiado el texto}
        Un sitio puede caber perfectamente y seguir siendo difícil de utilizar.

        \begin{itemize}
            \item Mantenga tamaños de lectura cómodos.
            \item Evite párrafos demasiado anchos.
            \item Use suficiente interlineado.
            \item Conserve contraste entre texto y fondo.
            \item No dependa únicamente del color para comunicar estados.
        \end{itemize}

        \Idea{
            Responsivo también significa conservar la legibilidad y la interacción,
            no únicamente evitar el desplazamiento horizontal.
        }
    \end{frame}

    \begin{frame}[fragile]{Botones y enlaces fáciles de utilizar}
        \begin{csscodedark}
            .accion {
                display: inline-flex;
                align-items: center;
                justify-content: center;
                min-height: 44px;
                padding: 0.75rem 1rem;
            }
        \end{csscodedark}

        En pantallas táctiles, los controles demasiado pequeños o muy cercanos
        dificultan la selección.

        \Practice{
            Revise no solo cómo se ve la página, sino también qué tan fácil
            resulta pulsar sus enlaces y botones.
        }
    \end{frame}

    \begin{frame}{Ocultar contenido requiere una razón}
        Cambiar de tamaño no significa eliminar automáticamente información.

        Antes de utilizar \texttt{display: none}, pregunte:

        \begin{itemize}
            \item ¿el contenido sigue siendo necesario?;
            \item ¿puede reorganizarse?;
            \item ¿puede resumirse sin perder significado?;
            \item ¿existe otra manera de acceder a él?
        \end{itemize}

        \Warning{
            La versión móvil no debe convertirse en una versión incompleta del sitio.
        }
    \end{frame}

    % =========================================================
    \section{Preferencias del usuario}
    % =========================================================

    \SectionSlide
    {Preferencias del usuario}
    {El diseño también responde a necesidades personales}

    \begin{frame}[fragile]{Reducir movimiento}
        Algunas personas prefieren menos animaciones:

        \begin{csscodedark}
            @media (prefers-reduced-motion: reduce) {
                *,
                *::before,
                *::after {
                    scroll-behavior: auto;
                    animation-duration: 0.01ms;
                    animation-iteration-count: 1;
                    transition-duration: 0.01ms;
                }
            }
        \end{csscodedark}

        \Idea{
            Las media queries también pueden responder a preferencias de accesibilidad,
            no solo al ancho de la ventana.
        }
    \end{frame}

    \begin{frame}[fragile]{Esquema de color preferido}
        Un sitio puede reconocer el modo oscuro del sistema:

        \begin{csscodedark}
            @media (prefers-color-scheme: dark) {
                body {
                    color: #eef7f2;
                    background-color: #10221c;
                }

                .tarjeta {
                    background-color: #19372d;
                }
            }
        \end{csscodedark}

        Esto es opcional y requiere verificar cuidadosamente el contraste.
    \end{frame}

    % =========================================================
    \section{Depuración}
    % =========================================================

    \SectionSlide
    {Probar y depurar}
    {Encontrar el ancho donde la interfaz deja de funcionar}

    \begin{frame}{Probar solo tres tamaños no es suficiente}
        No revise únicamente una vista móvil, una tableta y un escritorio.

        \begin{enumerate}
            \item Abra las herramientas de desarrollo.
            \item Active la vista adaptable.
            \item Cambie el ancho lentamente.
            \item Observe cada punto intermedio.
            \item Corrija donde el contenido realmente falla.
        \end{enumerate}

        \Practice{
            También cambie manualmente el tamaño de la ventana del navegador.
            Una página puede abrirse en una ventana estrecha dentro de un monitor grande.
        }
    \end{frame}

    \begin{frame}{Qué debemos observar}
        \begin{itemize}
            \item Desplazamiento horizontal inesperado.
            \item Texto cortado o demasiado estrecho.
            \item Imágenes deformadas.
            \item Controles que se superponen.
            \item Espacios excesivos o insuficientes.
            \item Cambios bruscos cerca de los breakpoints.
            \item Orden de lectura incoherente.
        \end{itemize}

        \Question[Comprobación]{
            ¿La página sigue siendo comprensible con zoom del navegador al 200\%?
        }
    \end{frame}

    \begin{frame}[fragile]{Encontrar el elemento que desborda}
        Durante la depuración puede aplicar temporalmente:

        \begin{csscodedark}
            * {
                outline: 1px solid red;
            }
        \end{csscodedark}

        Los contornos ayudan a reconocer qué elemento supera el ancho disponible.

        \Warning{
            Esta regla es temporal. Retírela cuando termine la revisión.
            No resuelva el problema ocultando todo el desbordamiento del documento.
        }
    \end{frame}

    \begin{frame}{Revise el HTML antes de culpar al CSS}
        Algunos problemas responsivos nacen en la estructura:

        \begin{itemize}
            \item imágenes sin contenedor adecuado;
            \item navegación sin agrupación semántica;
            \item elementos cerrados incorrectamente;
            \item contenido importante ubicado en un orden poco lógico;
            \item tablas utilizadas para construir la distribución.
        \end{itemize}

        \Idea{
            Un HTML ordenado facilita que la interfaz pueda reorganizarse sin perder sentido.
        }
    \end{frame}

    % =========================================================
    \section{Práctica guiada}
    % =========================================================

    \SectionSlide
    {Práctica guiada}
    {Biblioteca itinerante}

    \begin{frame}{Situación}
        Una biblioteca necesita una página para anunciar sus próximas visitas.

        La página debe incluir:

        \begin{itemize}
            \item encabezado e introducción;
            \item navegación interna;
            \item próximas paradas en tarjetas;
            \item una sección de inscripción;
            \item pie de página.
        \end{itemize}

        \Result{
            La información aparecerá en una columna cuando el espacio sea estrecho
            y se reorganizará progresivamente en anchos mayores.
        }
    \end{frame}

    \begin{frame}[fragile]{Estructura del proyecto}
        \begin{bashcodedark}
            biblioteca-itinerante/
            ├── index.html
            ├── css/
            │   └── styles.css
            └── img/
            ├── biblioteca.jpg
            ├── parque.jpg
            └── colegio.jpg
        \end{bashcodedark}

        \Practice{
            Construya primero la estructura y compruebe las rutas antes de aplicar estilos.
        }
    \end{frame}

    \begin{frame}[fragile]{Paso 1: configuración inicial}
        \begin{htmlcodedark}
        <head>
            <meta charset="UTF-8">
            <meta name="viewport"
                content="width=device-width, initial-scale=1.0">
            <title>Biblioteca itinerante</title>
            <link rel="stylesheet" href="css/styles.css">
        </head>
        \end{htmlcodedark}

        \begin{csscodedark}
            * { box-sizing: border-box; }
            body { margin: 0; font-family: Arial, sans-serif; }
            img { display: block; max-width: 100%; height: auto; }
        \end{csscodedark}
    \end{frame}

    \begin{frame}[fragile]{Paso 2: estructura principal}
        \begin{htmlcodedark}
        <header class="cabecera">
            <h1>Biblioteca itinerante</h1>
            <p>Los libros llegan a su comunidad.</p>
        </header>

        <nav class="menu">
            <a href="#paradas">Paradas</a>
            <a href="#inscripcion">Inscripción</a>
        </nav>

        <main class="contenido">
            <section class="presentacion">
                <img src="img/biblioteca.jpg" alt="Biblioteca móvil">
                <div><h2>Una biblioteca en movimiento</h2></div>
            </section>
        \end{htmlcodedark}
    \end{frame}

    \begin{frame}[fragile]{Paso 3: tarjetas}
        \begin{htmlcodedark}
        <section id="paradas" class="paradas">
            <article class="parada">
                <img src="img/parque.jpg" alt="Parque del barrio">
                <h3>Parque Central</h3>
                <p>Sábado, 9:00 a. m.</p>
            </article>

            <article class="parada">
                <img src="img/colegio.jpg" alt="Institución educativa">
                <h3>Colegio Horizonte</h3>
                <p>Miércoles, 2:00 p. m.</p>
            </article>
    </section>
    </main>
        \end{htmlcodedark}
    \end{frame}

    \begin{frame}[fragile]{Paso 4: versión para poco espacio}
        \vspace{-0.45cm}
        \begin{csscodedark}
            .contenido {
                width: min(92%, 1100px);
                margin-inline: auto;
            }

            .menu {
                display: flex;
                flex-direction: column;
                gap: 0.5rem;
                padding: 1rem;
            }

            .presentacion,
            .paradas {
                display: grid;
                gap: 1rem;
            }
        \end{csscodedark}
    \end{frame}

    \begin{frame}[fragile]{Paso 5: mejorar las tarjetas}
        \vspace{-0.45cm}
        \begin{csscodedark}
            .parada {
                overflow: hidden;
                border: 1px solid #b8c7bf;
                border-radius: 0.75rem;
                background-color: white;
            }

            .parada h3,
            .parada p {
                margin-inline: 1rem;
            }

            .parada img {
                width: 100%;
                height: 220px;
                object-fit: cover;
            }
        \end{csscodedark}
    \end{frame}

    \begin{frame}[fragile]{Paso 6: primer breakpoint}
        \begin{csscodedark}
            @media (min-width: 700px) {
                .menu {
                    flex-direction: row;
                    justify-content: center;
                }

                .presentacion {
                    grid-template-columns: 1.2fr 1fr;
                    align-items: center;
                }

                .paradas {
                    grid-template-columns: repeat(2, 1fr);
                }
            }
        \end{csscodedark}
    \end{frame}

    \begin{frame}[fragile]{Paso 7: ampliar el catálogo}
        Si se agregan más paradas, puede utilizar:

        \begin{csscodedark}
            @media (min-width: 1000px) {
                .paradas {
                    grid-template-columns: repeat(3, 1fr);
                }
            }
        \end{csscodedark}

        O reemplazar los breakpoints de las tarjetas por:

        \begin{csscodedark}
            .paradas {
                grid-template-columns:
                repeat(auto-fit, minmax(min(100%, 250px), 1fr));
            }
        \end{csscodedark}
    \end{frame}

    \begin{frame}{Verificación de la práctica}
        Revise la página en los siguientes escenarios:

        \begin{itemize}
            \item 360px: contenido en una columna y sin desbordamiento.
            \item 699px: comportamiento inmediatamente antes del breakpoint.
            \item 700px: aplicación correcta de la nueva composición.
            \item 999px y 1000px: transición hacia tres columnas, si existe.
            \item Ventana libre: arrastre el borde lentamente.
            \item Zoom al 200\%: el contenido continúa disponible.
        \end{itemize}

        \Question[Explique]{
            ¿Qué reglas son base y cuáles solo se aplican cuando existe más espacio?
        }
    \end{frame}

    % =========================================================
    \section{Ejercicio autónomo}
    % =========================================================

    \SectionSlide
    {Ejercicio autónomo}
    {Guía de mercados locales}

    \begin{frame}{Situación propuesta}
        Construya una página para presentar mercados y emprendimientos locales.

        Debe incluir como mínimo:

        \begin{itemize}
            \item encabezado y navegación;
            \item una presentación principal con imagen;
            \item cuatro tarjetas de mercados o emprendimientos;
            \item horario, ubicación y descripción breve;
            \item un llamado a la acción;
            \item pie de página.
        \end{itemize}

        El ejercicio es individual.
    \end{frame}

    \begin{frame}{Requisitos técnicos}
        Utilice como mínimo:

        \begin{itemize}
            \item etiqueta \texttt{meta viewport};
            \item imágenes con \texttt{max-width: 100\%};
            \item un contenedor con ancho flexible y límite máximo;
            \item Flexbox o Grid en la navegación;
            \item Grid para organizar las tarjetas;
            \item enfoque mobile first;
            \item al menos dos media queries con \texttt{min-width};
            \item textos alternativos en las imágenes.
        \end{itemize}
    \end{frame}

    \begin{frame}[fragile]{Estructura de entrega}
        \begin{bashcodedark}
            mercados-nombre-apellido/
            ├── index.html
            ├── css/
            │   └── styles.css
            └── img/
            ├── portada.jpg
            ├── mercado-1.jpg
            ├── mercado-2.jpg
            ├── mercado-3.jpg
            └── mercado-4.jpg
        \end{bashcodedark}

        Todos los archivos deben abrirse correctamente desde la carpeta entregada.
    \end{frame}

    \begin{frame}{Orden recomendado de trabajo}
        \begin{enumerate}
            \item Construya un HTML semántico y ordenado.
            \item Conecte la hoja de estilos y compruebe las rutas.
            \item Termine primero la versión de menor ancho.
            \item Identifique dónde el contenido necesita reorganizarse.
            \item Agregue las media queries.
            \item Pruebe anchos intermedios y corrija desbordamientos.
            \item Añada mejoras visuales únicamente cuando lo básico funcione.
        \end{enumerate}

        \Warning{
            Los elementos adicionales no reemplazan ninguno de los requisitos mínimos.
        }
    \end{frame}

    \begin{frame}{Lista de comprobación}
        Antes de entregar, confirme:

        \begin{itemize}
            \item No existe desplazamiento horizontal accidental.
            \item El contenido se comprende en una sola columna.
            \item Las imágenes no se deforman.
            \item La navegación cambia correctamente.
            \item Las tarjetas aprovechan el espacio disponible.
            \item Los breakpoints responden a necesidades del contenido.
            \item El texto sigue siendo legible.
            \item Puede explicar cada regla utilizada.
        \end{itemize}
    \end{frame}

    % =========================================================
    \section{Errores frecuentes}
    % =========================================================

    \SectionSlide
    {Errores frecuentes}
    {Problemas que impiden una adaptación correcta}

    \begin{frame}{Diseñar únicamente para tamaños conocidos}
        Una interfaz puede verse bien a 375px y 1024px, pero fallar a 620px.

        \Warning{
            Los tamaños predeterminados de las herramientas del navegador
            son puntos de prueba, no una garantía de funcionamiento responsivo.
        }

        Cambie el ancho de manera continua para descubrir los estados intermedios.
    \end{frame}

    \begin{frame}{Crear demasiados breakpoints}
        No agregue una media query cada vez que cambia unos pocos píxeles.

        Esto produce:

        \begin{itemize}
            \item reglas difíciles de mantener;
            \item sobrescrituras contradictorias;
            \item comportamiento impredecible;
            \item dependencia de dispositivos específicos.
        \end{itemize}

        \Idea{
            Comience con una estructura flexible y añada un breakpoint
            solamente cuando el contenido lo necesite.
        }
    \end{frame}

    \begin{frame}{Usar anchos fijos en toda la página}
        Una declaración como \texttt{width: 1200px} puede funcionar en una pantalla amplia,
        pero no tiene espacio suficiente dentro de un viewport de 360px.

        Prefiera combinaciones como:

        \begin{center}
            \texttt{width: 92\%; max-width: 1200px;}
        \end{center}

        Los tamaños fijos siguen siendo útiles para iconos, bordes o límites concretos.
    \end{frame}

    \begin{frame}{Confundir responsivo con ocultar}
        Ocultar navegación, texto o acciones puede hacer que la interfaz parezca limpia,
        pero también puede impedir que el usuario complete una tarea.

        \Compare{
            Reorganizar conserva el contenido y cambia su presentación.
            Ocultar elimina su disponibilidad en ese contexto.
        }
    \end{frame}

    \begin{frame}{Cambiar el orden visual sin revisar la lectura}
        Grid y Flexbox pueden modificar el orden visual de los elementos.

        Sin embargo:

        \begin{itemize}
            \item el teclado normalmente sigue el orden del HTML;
            \item los lectores de pantalla interpretan la estructura del documento;
            \item el orden semántico debe conservar sentido sin CSS.
        \end{itemize}

        \Warning{
            No utilice \texttt{order} para corregir una estructura HTML mal organizada.
        }
    \end{frame}

    \begin{frame}{Corregir todo con overflow hidden}
        Aplicar \texttt{overflow-x: hidden} al documento puede ocultar la evidencia,
        pero no corrige el elemento que causa el desbordamiento.

        \Practice{
            Active contornos temporales, inspeccione tamaños calculados
            y localice el elemento responsable.
        }

        La corrección puede requerir \texttt{max-width}, \texttt{min-width: 0},
        ajuste de línea o una distribución diferente.
    \end{frame}

    % =========================================================
    \section{Cierre}
    % =========================================================

    \SectionSlide
    {Cierre}
    {Decidir qué debe fluir y qué debe cambiar}

    \begin{frame}{Mapa de decisiones}
        \begin{center}
            \begin{tikzpicture}[
                scale=0.88,
                transform shape,
                nodo/.style={draw=BrandBlue, thick, rounded corners,
                    minimum width=3.4cm, minimum height=0.85cm, align=center},
                flecha/.style={->, >=Latex, thick}
                ]
                \node[nodo, fill=IdeaBody] (p) at (5.5,5.3) {¿El contenido se desborda?};
                \node[nodo, fill=PracticeBody] (f) at (2.2,2.6) {Use medidas flexibles\\e imágenes adaptables};
                \node[nodo, fill=PracticeBody] (c) at (8.8,2.6) {¿La composición deja\\de ser clara?};
                \node[nodo, fill=ResultBody] (m) at (8.8,0.8) {Agregue una media query};
                \draw[flecha] (p) -- (f) node[midway, above left] {sí};
                \draw[flecha] (p) -- (c) node[midway, above right] {no};
                \draw[flecha] (c) -- (m) node[midway, right] {sí};
            \end{tikzpicture}
        \end{center}

        Primero permita que la interfaz fluya; después cambie la composición
        solo cuando resulte necesario.
    \end{frame}

    \begin{frame}{Consulta rápida}
        \begin{center}
            \small
            \begin{tabular}{p{5.0cm}p{7.5cm}}
                \toprule
                \textbf{Recurso} & \textbf{Propósito} \\
                \midrule
                \texttt{meta viewport} & usar el ancho real del viewport \\
                \texttt{max-width: 100\%} & limitar imágenes y otros elementos \\
                \texttt{max-width} & impedir un crecimiento excesivo \\
                \texttt{@media (min-width: ...)} & mejorar la composición al ganar espacio \\
                \texttt{flex-wrap} & permitir varias líneas en Flexbox \\
                \texttt{auto-fit + minmax()} & crear columnas automáticas en Grid \\
                \texttt{clamp()} & establecer tamaño mínimo, flexible y máximo \\
                \bottomrule
            \end{tabular}
        \end{center}
    \end{frame}

    \begin{frame}{Preguntas de comprensión}
        \begin{enumerate}
            \item ¿Qué diferencia existe entre viewport y pantalla física?
            \item ¿Por qué \texttt{max-width} ayuda a construir interfaces fluidas?
            \item ¿Cuándo se justifica crear un breakpoint?
            \item ¿Cómo se interpreta \texttt{min-width: 700px}?
            \item ¿Qué caracteriza al enfoque mobile first?
            \item ¿Cuándo puede Grid adaptarse sin una media query?
            \item ¿Por qué no conviene bloquear el zoom?
        \end{enumerate}
    \end{frame}

    \begin{frame}{Lo esencial de esta clase}
        \begin{itemize}
            \item El diseño responsivo se basa en el espacio disponible.
            \item El viewport debe configurarse desde el HTML.
            \item Las medidas flexibles resuelven gran parte de la adaptación.
            \item Las imágenes deben respetar el ancho de su contenedor.
            \item Los breakpoints nacen de las necesidades del contenido.
            \item Mobile first construye desde una base sencilla hacia espacios mayores.
            \item Flexbox y Grid permiten reorganizar la interfaz.
            \item La verificación debe incluir anchos intermedios, zoom e interacción.
        \end{itemize}

        \Result{
            Una página responsiva no intenta adivinar el dispositivo:
            conserva su utilidad mientras cambia el espacio disponible.
        }
    \end{frame}

    \begin{frame}{Aplicación en el proyecto intermedio}
        En el proyecto del concierto deberán comprobar que:

        \begin{itemize}
            \item las dos páginas funcionan en anchos estrechos y amplios;
            \item cada hoja CSS conserva una organización clara;
            \item la navegación permite pasar entre las páginas;
            \item las imágenes no generan desbordamiento;
            \item Flexbox, Grid y las media queries se utilizan con una finalidad explicable.
        \end{itemize}

        \Warning{
            Primero aseguren el cumplimiento de lo solicitado.
            Después pueden incorporar elementos adicionales sin sustituir los requisitos básicos.
        }
    \end{frame}

\end{document}
